\section{Formal-group height and the next counting invariant}\label{fg:sec:height}
The coherent line also detects the first height of an Artin--Mazur formal
group. Higher heights require integral coefficient data of greater
$p$-adic precision. We first establish the geometric interpretation and
the exact precision needed, before analyzing the counting construction.

\subsection{The top formal group and its first coefficient}
Let $k=\Fp$, let $d\ge3$, and let $X=V(F)\subset\PP^{d-1}_k$ be a smooth
geometrically integral hypersurface defined by a homogeneous polynomial
of degree $d$. Write $n=d-2$. For a local Artinian $k$-algebra $A$ with
residue field $k$, define
\[
 \Phi_X(A)=\ker\bigl(H^n_{\mathrm{et}}(X_A,\mathbb G_m)
                    \longrightarrow H^n_{\mathrm{et}}(X,\mathbb G_m)\bigr).
\]
Here $X_A=X\times_k\Spec A$. The height of a one-dimensional formal group
in characteristic $p$ is $h$ if the first nonzero term of its
multiplication-by-$p$ series has degree $p^h$; it is infinite if that series
is zero. Thus $h\ge r$ always includes infinite height.

\begin{proposition}[The coherent line and height one]\label{fg:prop:height-one}
The functor $\Phi_X$ is a smooth one-dimensional commutative formal group.
In a coordinate with tangent generator as in Lemma~\ref{lem:coherent},
\[
 [p]_{\Phi_X}(T)=H_p(F)T^p\pmod{T^{p+1}},\qquad
 H_p(F)=[(x_1\cdots x_d)^{p-1}]F^{p-1}.
\]
Consequently
\[
 h(X)=1\iff H_p(F)\ne0\iff\#X(\Fp)\not\equiv1\pmod p.
\]
On the family of Theorem~\ref{thm:main}, this scalar is $S\bmod p$.
For fixed $p$, both $h(X)=1$ and $h(X)\ge2$ are $\ModP$-complete
promise problems for circuit input in growing dimension. Each is in
deterministic polynomial time for explicit sparse input, or for circuit
input in fixed dimension. In dimension one the only heights are one and
two.
\end{proposition}

\begin{proof}
The hypersurface sequence gives
\[
 H^0(X,\OO_X)=k,\quad H^i(X,\OO_X)=0\ (0<i<n),\quad
 H^n(X,\OO_X)=k,\quad H^{n+1}(X,\OO_X)=0.
\]
For $n\ge2$, pro-representability follows from
\citet[II, Corollary~4.2, p.~109]{fg-artin-mazur}; the next-degree
vanishing gives smoothness, and the tangent space is $H^n(X,\OO_X)$
\citep[II, Corollaries~2.4--2.5 and~2.11--2.12,
pp.~106--107]{fg-artin-mazur}. Their complete-intersection special case
on p.~110 gives the same conclusion.

For $n=1$, the previous-degree vanishing would be false. Instead the
degree-zero deformation functor is $A\mapsto1+\mathfrak m_A$, since
$H^0(X_A,\OO_{X_A})=A$. It is formally smooth. Corollary~2.12 of
Artin--Mazur therefore applies again, since $H^2(X,\OO_X)=0$.
Their identification in II, \S4, p.~109, makes $\Phi_X$ the formal
completion of the Picard functor at the identity. After extension to an
algebraic closure, choosing a point identifies this formal group with
that of the resulting elliptic curve. Its height is one or two
\citep[IV, Theorem~7.4 and Corollary~7.5, p.~134]{fg-silverman}.
Height is unchanged by a field extension, so this holds over $k$ as well.

We verify the first coefficient directly, including its normalization.
For $I=\mathfrak m_A$, the constant family splits the reduction map on
unit sheaves. Thus
\[
 \Phi_X(A)=H^n(X,1+I\OO_X).
\]
The successive quotients of the $I$-adic filtration are coherent additive
sheaves. They are acyclic on affine opens, so a finite affine cover
computes this cohomology by multiplicative \v Cech cocycles; it also
identifies its Zariski and \'etale descriptions. This is the
principal-unit description of \citet[II.(1.3)--(1.4), pp.~103--104,
and Corollary~2.6, p.~106]{fg-artin-mazur}.

Choose a nonzero tangent vector $\eta\in H^n(X,\OO_X)$ and the matching
formal coordinate. Over $k[t]/(t^{p+1})$, the point with coordinate $t$
has a cocycle representative with components $1+t a_I+O(t^2)$, where
$(a_I)$ represents $\eta$. Multiplication by $p$ raises each component
to its $p$-th power, giving
\[
 (1+t a_I+O(t^2))^p=1+t^p a_I^p\pmod{t^{p+1}}.
\]
If coherent Frobenius carries $\eta$ to $c\eta$, naturality of the tangent
identification under $k[\epsilon]/(\epsilon^2)\to k[t]/(t^{p+1})$,
$\epsilon\mapsto t^p$, identifies this point with coordinate $ct^p$.
Hence $[p](T)=cT^p\bmod T^{p+1}$. Lemma~\ref{lem:coherent} gives
$c=H_p(F)$, and equation~\eqref{eq:point-sign} gives the point-count
equivalence. Over $\Fp$ this scalar is unchanged by a nonzero change of
tangent basis: the semilinear rescaling factor is $\lambda^{p-1}=1$.

For completeness, the first nonzero degree of a homomorphism of formal
groups in characteristic $p$ is a power of $p$. Its leading term
$aT^e$ satisfies $(X+Y)^e=X^e+Y^e$ by comparing total degree $e$ in the
homomorphism identity. The base-$p$ expansion of $(1+X)^e$ forces $e$ to
be a $p$-power. The linear term of $[p]$ is zero. This proves the claimed
height criterion, including the zero-series case.

A nondeterministic machine counts the projective rational points by
guessing the unique representative whose first nonzero coordinate is
one. Adding $p-1$ accepting paths gives a count congruent to
$\#X(\Fp)-1$. Smoothness forces every ambient variable to occur in the
circuit, since an omitted variable produces a singular coordinate
vertex. The number of guessed coordinates is therefore polynomial in
the input length, even when their labels are binary encoded.
This proves membership for height one. If $a$ is any accepting-path
count, then
\[
 1+(p-1)a^{p-1}\equiv
 \begin{cases}1,&a\equiv0\pmod p,\\0,&a\not\equiv0\pmod p.
 \end{cases}
\]
Sums and fixed powers are again accepting-path counts. Thus $\ModP$ is
closed under complement, with this explicit fixed-prime construction.
Theorem~\ref{thm:main} and parsimonious width-three encodings give
hardness for both predicates. The same reduction permits any fixed
lower bound on the output dimension by choosing a larger $p$-power $Q$.
For sparse input, enumerate the $t^{p-1}$ ordered term selections in
$F^{p-1}$. In fixed dimension, enumerate the fixed finite set of
projective $\Fp$-points. These give the asserted polynomial algorithms.
\end{proof}

\subsection{Integral logarithms and every height test}
Let now $k$ be any perfect field of characteristic $p$, put $R=W(k)$,
and normalize $v_p(p)=1$. Choose a homogeneous lift
$\widetilde F\in R[x_1,\ldots,x_d]$ of the equation, and define
\[
 b_m=[(x_1\cdots x_d)^{m-1}]\widetilde F^{m-1},\qquad b_1=1.
\]
\citet[Theorem~1, p.~907; specialization on p.~908 and in
\S3.9, p.~914]{fg-stienstra} constructs an integral formal-group law
$\mathcal F$ lifting $\Phi_X$, with normalized logarithm
\begin{equation}\label{fg:eq:stienstra}
 L(T)=\sum_{m\ge1}\frac{b_m}{m}T^m.
\end{equation}
Indeed, the positive integer vectors of length $d$ and total $d$ consist
only of $(1,\ldots,1)$, so his index set has one element. The lift is
flat: if $pA=\widetilde F B$, reduction modulo $p$ in the polynomial
domain gives $B=pB_1$, and cancellation gives $A=\widetilde F B_1$.
Thus the quotient is torsion-free over the discrete valuation ring.
The defining sequence has one nonzero element and is regular. These
verify the base and flatness hypotheses in Stienstra's theorem.
The coefficient of $T^p$ in \eqref{fg:eq:stienstra} is $b_p/p$;
the coherent scalar is the reduction of its integral numerator $b_p$.

\begin{lemma}[Logarithmic test for the next height]\label{fg:lem:height-log}
Let $\mathcal F$ be any one-dimensional commutative formal-group law
over $R$, with normalized logarithm
$L(T)=T+\sum_{m\ge2}\ell_mT^m$. If its reduction has height at least
$r\ge1$, and $q=p^r$, then
\[
 \ell_m\in R\ (m<q),\qquad p\ell_q\in R,
 \qquad p\ell_q\equiv[T^q][p]_{\mathcal F}(T)\pmod p.
\]
Its height is at least $r+1$ precisely when $p\ell_q\in pR$.
For the logarithm \eqref{fg:eq:stienstra}, this says
\begin{equation}\label{fg:eq:next-height}
 b_{p^r}\in p^{r-1}R,\qquad
 h(X)\ge r+1\iff b_{p^r}\in p^rR
 \quad\text{if }h(X)\ge r.
\end{equation}
In particular,
\begin{align}
 h(X)\ge r&\iff p^j\mid b_{p^j}\quad(1\le j<r),\label{fg:eq:all-height}\\
 h(X)=r&\iff p^j\mid b_{p^j}\quad(1\le j<r),
                  \quad p^r\nmid b_{p^r},\label{fg:eq:exact-height}\\
 h(X)=\infty&\iff p^j\mid b_{p^j}\quad\text{for every }j\ge1.
 \label{fg:eq:infinite-height}
\end{align}
\end{lemma}

\begin{proof}
Write $g(T)=[p]_{\mathcal F}(T)=pT+\sum_{m\ge2}\gamma_mT^m$.
Comparing coefficients in $L(g(T))=pL(T)$ gives
\begin{equation}\label{fg:eq:log-recursion}
 (p-p^m)\ell_m=\gamma_m+
       \sum_{i=2}^{m-1}\ell_i[T^m]g(T)^i.
\end{equation}
The height hypothesis means $\gamma_m\in pR$ for $m<q$.
Inductively, every summand on the right for $m<q$ is in $pR$,
and $(p-p^m)/p$ is a unit. Thus $\ell_m$ is integral there.
At $m=q$, each coefficient in $[T^q]g^i$ for $i\ge2$ uses only
degrees smaller than $q$ and is divisible by $p^i$. Hence
$p\ell_q$ is integral and reduces to $\gamma_q$. The first possible
nonzero term of the reduced $p$-series has $p$-power degree, as proved
above. Its coefficient at $q$ vanishes exactly when the height is at
least $r+1$. Substituting $\ell_q=b_q/q$ and induction in $r$ proves
all the assertions, including the infinite case.
\end{proof}

This proof uses the original coordinate and does not discard the
non-$p$-power terms of its logarithm. For comparison, an integral
$p$-typicalization keeps the coefficients
$\lambda_i=b_{p^i}/p^i$. By
\citet[Theorem~2.8 and (4.3.1), pp.~134,138]{fg-hazewinkel},
its Hazewinkel parameters satisfy
\[
 v_n=p\lambda_n-\sum_{i=1}^{n-1}\lambda_i v_{n-i}^{p^i}\in R.
\]
In particular
\[
 v_1=b_p,\qquad v_2=\frac{b_{p^2}-b_p^{p+1}}p.
\]
They are parameters of the integral $p$-typical law, rather than all the
ordinary power-series coefficients of $[p]$ in the original coordinate.
The first nonzero one modulo $p$ detects the same height. The recursion
above is not needed to infer \eqref{fg:eq:next-height}.

\begin{corollary}[The second height invariant]\label{fg:cor:second}
On the locus $H_p(F)=0$, the element $b_{p^2}$ is divisible by $p$, and
\[
 h(X)\ge3\iff\frac{b_{p^2}}p=0\pmod p
            \iff b_{p^2}=0\pmod{p^2}.
\]
If $a_{p^2}$ denotes the literal coefficient of $T^{p^2}$ in $L$, this
invariant is $p a_{p^2}\bmod p$. The coefficient $a_{p^2}$ itself need
not be integral. For a fixed defining polynomial over $k$, the residue
$b_{p^2}/p\bmod p$ on this locus is independent of its homogeneous lift.
\end{corollary}

\begin{proof}
Only lift independence remains to be checked. Replace
$\widetilde F$ by $\widetilde F+pH$, with $H$ homogeneous of degree $d$.
Modulo $p^2$, the difference of the numerators is
\[
 p(p^2-1)[(x_1\cdots x_d)^{p^2-1}]H\widetilde F^{p^2-2}.
\]
Modulo $p$, the product after the coefficient sign is
$H F^{p-2}(F^{p-1})^p$. A contributing exponent of $HF^{p-2}$ is
congruent to $(p-1,\ldots,p-1)$ coordinatewise modulo $p$.
Its total degree is $d(p-1)$, so nonnegativity forces equality with
that vector. The coefficient is therefore
\[
 [(x_1\cdots x_d)^{p-1}]HF^{p-2}\, H_p(F)^p=0.
\]
This proves equality of the numerators modulo $p^2$, and hence of the
quotient residues. General height tests are intrinsic by their
equivalence to the reduced $p$-series. Under a nonstrict rescaling
$U=uT$, its first nonzero coefficient has weight $u^{1-p^h}$, so it is
its vanishing that is independent of every tangent normalization.
\end{proof}

\begin{example}[Why the second numerator needs two digits]\label{fg:ex:fermat}
For the Fermat cubic over $\F_2$, the integer lift gives $b_2=0$ and
$b_4=3!=6$. Thus $b_4/2=1\bmod2$ and its height is exactly two.
Both numerators vanish modulo two, so their bare reductions cannot
distinguish the two stages. Replacing the lift by $F+2xyz$ gives
$b'_2=2$ and $b'_4=14$, with the same second residue; the unrestricted
parameter is $(14-2^3)/2=3$.

For the Fermat quartic over $\F_3$, $b_3=0$ and
$b_9=8!/(2!)^4=2520$, divisible by nine. Its height is infinite.
Indeed, $b_{3^r}=0$ for odd $r$. For $r=2s$ put
$a=(3^{2s}-1)/4=2(1+9+\cdots+9^{s-1})$. Its base-three digit sum is
$2s=r$, and that of $3^r-1$ is $2r$. Legendre's factorial formula gives
\[
 v_3\left(\frac{(3^r-1)!}{(a!)^4}\right)=r.
\]
Equation~\eqref{fg:eq:infinite-height} applies at every stage.
These are smooth equations, since their coordinate derivatives have
no common projective zero. The multinomial formula also appears in
\citet[Example~4.13, p.~919]{fg-stienstra}.
\end{example}
